\chapter{1977年广西百色卷（理）}

\section{解答题}

\begin{problem}
求函数 \(y=\dfrac{2\sqrt{5-x}}{\sqrt{x-3}}\) 的定义域.
\end{problem}

\begin{solution}
要使分子中的根式有意义，必须有 \(5-x\geqslant0\)，所以 \(x\leqslant5\). 分母中的根式既要有意义又不能为零，因此必须有 \(x-3>0\)，所以 \(x>3\). 两个条件合并得 \(3<x\leqslant5\)，故函数的定义域为 \(\left(3,5\right]\).
\end{solution}

\begin{problem}
计算 \(\left(-\dfrac{1}{2}\right)^{-3}+32^{\frac{2}{5}}+5^0-3^{-2}\).
\end{problem}

\begin{solution}
负整数指数表示相应正整数次幂的倒数，且 \(32=2^5\)，所以
\[
\left(-\dfrac{1}{2}\right)^{-3}+32^{\frac{2}{5}}+5^0-3^{-2}=(-2)^3+\left(2^5\right)^{\frac{2}{5}}+1-\dfrac{1}{3^2}=-8+4+1-\dfrac{1}{9}=-\dfrac{28}{9}
\]
\end{solution}

\begin{problem}
求 \(\sin^2 60^\circ+\cot45^\circ+\dfrac{3}{4}\tan^2 30^\circ-\cos60^\circ\) 的值.
\end{problem}

\begin{solution}
由特殊角三角函数值 \(\sin60^\circ=\dfrac{\sqrt3}{2}\)，\(\cot45^\circ=1\)，\(\tan30^\circ=\dfrac{\sqrt3}{3}\)，\(\cos60^\circ=\dfrac12\)，原式为
\[
\left(\dfrac{\sqrt3}{2}\right)^2+1+\dfrac34\left(\dfrac{\sqrt3}{3}\right)^2-\dfrac12=\dfrac34+1+\dfrac14-\dfrac12=\dfrac32
\]
\end{solution}

\begin{problem}
\(\log_2\sqrt8=x\). 求 \(x\).
\end{problem}

\begin{solution}
由对数的定义，\(\log_2\sqrt8=x\) 等价于 \(2^x=\sqrt8\). 又因为 \(\sqrt8=\sqrt{2^3}=2^{\frac32}\)，所以 \(2^x=2^{\frac32}\). 底数 \(2\) 大于 \(0\) 且不等于 \(1\)，幂函数的值相等时指数相等，故 \(x=\dfrac32\).
\end{solution}

\begin{problem}
已知直线过 \((3,1)\) 和 \((-6,4)\) 两点.

\begin{enumerate}
    \item 求出这直线方程.
    \item 求它的斜率和 \(y\) 轴上的截距.
    \item 这函数是上升或下降.
    \item 求过 \((0,-5)\) 和这直线平行的直线方程.
\end{enumerate}
\end{problem}

\begin{solution}
\begin{enumerate}
\item 设直线方程为 \(y=kx+b\). 将两点坐标代入，得到
\[
\begin{cases}
1=3k+b,\\
4=-6k+b.
\end{cases}
\]
两式相减得 \(3=-9k\)，所以 \(k=-\dfrac13\). 代回 \(1=3k+b\)，得 \(b=2\). 因而直线方程为 \(y=-\dfrac13x+2\)，化为一般式为 \(x+3y-6=0\).

\item 由直线的斜截式 \(y=kx+b\) 可知，斜率为 \(k=-\dfrac13\)，\(y\) 轴上的截距为 \(b=2\).

\item 因为斜率 \(k=-\dfrac13<0\)，所以函数 \(y=-\dfrac13x+2\) 随 \(x\) 增大而减小，是下降函数.

\item 平行直线的斜率相等，所以所求直线的斜率仍为 \(-\dfrac13\). 设其方程为 \(y=-\dfrac13x+b_1\)，代入点 \((0,-5)\) 得 \(-5=b_1\). 因而所求直线为 \(y=-\dfrac13x-5\)，化为一般式为 \(x+3y+15=0\).
\end{enumerate}
\end{solution}

\begin{problem}
已知圆 \(O\) 中两弦 \(AB\)，\(CD\) 相交于 \(E\).

\begin{enumerate}
    \item 证明 \(AE\cdot BE=CE\cdot DE\).
    \item 当 \(CD\) 为直径 \(2R\)，\(AB\perp CD\) 时，证明 \(BE^2=CE(2R-CE)\).
\end{enumerate}

\bitmapfigure[max width=0.34\linewidth]{img/1977/guangxi_baise_science/q6_fig1.png}
\end{problem}

\begin{solution}
\begin{enumerate}
\item 连接 \(AC\)、\(BD\). 因为 \(A\)、\(E\)、\(B\) 共线，\(C\)、\(E\)、\(D\) 共线，所以 \(\angle AEC=\angle DEB\) 是对顶角. 又因为 \(A\)、\(B\)、\(C\)、\(D\) 在同一圆上，\(\angle ACE=\angle ACD\) 与 \(\angle DBE=\angle DBA\) 都对着弦 \(AD\)，故 \(\angle ACE=\angle DBE\). 因此 \(\triangle AEC\sim\triangle DEB\)，从而
\[
\dfrac{AE}{DE}=\dfrac{CE}{BE}
\]
交叉相乘得 \(AE\cdot BE=CE\cdot DE\).

\item 因为 \(CD\) 是过圆心 \(O\) 的直径，且 \(AB\perp CD\)，所以过圆心的垂线平分弦 \(AB\). 交点 \(E\) 是弦 \(AB\) 的中点，故 \(AE=BE\). 又 \(C\)、\(E\)、\(D\) 依次在直径上，且 \(CD=2R\)，所以 \(DE=CD-CE=2R-CE\). 将这两个关系代入第一问的结论，得
\[
BE^2=AE\cdot BE=CE\cdot DE=CE(2R-CE)
\]
\end{enumerate}
\end{solution}

\begin{problem}
一船往返于 \(48\text{ 公里}\) 的两地，共航行 \(10\text{ 小时}\)，已知水流速度是 \(2\text{ 公里/小时}\)，求这轮船在静水中航行的速度.
\end{problem}

\begin{solution}
设轮船在静水中的速度为 \(x\text{ 公里/小时}\). 轮船逆水航行时速度为 \(x-2\text{ 公里/小时}\)，必须有 \(x>2\)；顺水航行时速度为 \(x+2\text{ 公里/小时}\). 因此往返所用时间满足
\[
\dfrac{48}{x+2}+\dfrac{48}{x-2}=10
\]
由 \(x>2\) 可知 \((x+2)(x-2)\ne0\)，两边同乘该因式，得
\[
48(x-2)+48(x+2)=10(x+2)(x-2)
\]
整理得 \(5x^2-48x-20=0\). 因为
\[
5x^2-48x-20=(5x+2)(x-10)
\]
所以 \(x=10\) 或 \(x=-\dfrac25\). 其中 \(x=-\dfrac25\) 不满足速度必须为正且 \(x>2\) 的实际条件，舍去；\(x=10\) 代回原方程成立. 故轮船在静水中的速度为 \(10\text{ 公里/小时}\).
\end{solution}

\begin{problem}
已知 \(A\)，\(B\)，\(C\) 是 \(\triangle ABC\) 的三个内角，求证 \(a=b\cos C+c\cos B\)，\(b=c\cos A+a\cos C\)，\(c=a\cos B+b\cos A\).

\bitmapfigure[max width=0.34\linewidth]{img/1977/guangxi_baise_science/q8_fig1.png}
\end{problem}

\begin{solution}
作高 \(AD\perp BC\)，必要时将边延长. 在直角三角形 \(ABD\) 和 \(ACD\) 中，边 \(AB=c\) 在 \(BC\) 方向上的投影为 \(c\cos B\)，边 \(AC=b\) 在 \(BC\) 方向上的投影为 \(b\cos C\). 因此
\[
a=BC=c\cos B+b\cos C=b\cos C+c\cos B
\]

作高 \(BE\perp AC\)，必要时将边延长. 边 \(AB=c\) 在 \(AC\) 方向上的投影为 \(c\cos A\)，边 \(BC=a\) 在 \(AC\) 方向上的投影为 \(a\cos C\)，所以
\[
b=AC=c\cos A+a\cos C
\]

作高 \(CF\perp AB\)，必要时将边延长. 边 \(AC=b\) 在 \(AB\) 方向上的投影为 \(b\cos A\)，边 \(BC=a\) 在 \(AB\) 方向上的投影为 \(a\cos B\)，所以
\[
c=AB=b\cos A+a\cos B=a\cos B+b\cos A
\]
\end{solution}

\begin{problem}
当 \(a\)，\(b\)，\(c\) 为实数时，求证方程 \(x^2-(a+b)x+(ab-c^2)=0\) 有两个实数根，并求出两个根相等的条件.
\end{problem}

\begin{solution}
该方程是一元二次方程，二次项系数为 \(1\). 它的判别式为
\[
\Delta=\left[-(a+b)\right]^2-4\left(ab-c^2\right)=a^2+2ab+b^2-4ab+4c^2=(a-b)^2+4c^2
\]
因为 \(a\)、\(b\)、\(c\) 都是实数，\((a-b)^2\geqslant0\) 且 \(4c^2\geqslant0\)，所以 \(\Delta\geqslant0\). 由一元二次方程根的判别式判定，该方程有两个实数根（重根按两个根计）.

两个根相等当且仅当 \(\Delta=0\). 两个非负数之和为零，当且仅当它们都为零，因此 \((a-b)^2=0\)，\(4c^2=0\)，这等价于 \(a=b\)，\(c=0\). 此时重根为 \(x=\dfrac{a+b}{2}=a\).
\end{solution}

\section{选作题}

\begin{problem}
已知 \(2\lg(x-2y)=\lg x+\lg y\). 求 \(x:y\).
\end{problem}

\begin{solution}
原方程有意义的必要条件是 \(x>0\)，\(y>0\)，\(x-2y>0\). 在这些条件下，利用对数的运算性质，原方程可化为
\[
\lg\left(x-2y\right)^2=\lg x+\lg y=\lg(xy)
\]
因为常用对数函数在正数上是单调函数，所以
\[
\left(x-2y\right)^2=xy
\]
展开并移项，得 \(x^2-5xy+4y^2=0\)，即
\[
(x-y)(x-4y)=0
\]
故可能有 \(x=y\) 或 \(x=4y\). 但由原方程的定义域条件 \(x-2y>0\)，可得 \(\dfrac{x}{y}>2\)，所以 \(x=y\) 不符合条件；而 \(x=4y\) 满足 \(x>0\)、\(y>0\) 且 \(x-2y=2y>0\). 因此 \(\dfrac{x}{y}=4\)，即 \(x:y=4:1\).
\end{solution}

\begin{problem}
求内接于半径为 \(10\text{ 厘米}\)，圆心角为 \(120^\circ\) 的弓形的正方形的边长.

\bitmapfigure[max width=0.42\linewidth]{img/1977/guangxi_baise_science/q11_fig1.png}
\end{problem}

\begin{solution}
设图中正方形为 \(CDEF\)，其边长为 \(x\text{ 厘米}\). 作半径 \(OP\perp FE\)，令 \(N=OP\cap FE\)，\(M=OP\cap AB\). 由于过圆心的半径垂直于弦 \(FE\)，所以 \(N\) 是弦 \(FE\) 的中点；又因为 \(OP\perp AB\)，所以 \(M\) 是弦 \(AB\) 的中点，且 \(OP\) 平分圆心角 \(\angle AOB\). 因此 \(\angle BOM=\dfrac12\angle AOB=60^\circ\).

在直角三角形 \(OMB\) 中，\(OB=10\text{ 厘米}\)，所以
\[
OM=OB\cos60^\circ=10\times\dfrac12=5\text{ 厘米}
\]
正方形的高为边长，故 \(MN=x\)，而 \(N\) 是正方形上边 \(FE\) 的中点，故 \(NE=\dfrac{x}{2}\). 又 \(ON=OM+MN=5+x\). 在直角三角形 \(ONE\) 中，斜边 \(OE\) 是圆的半径，故 \(OE=10\text{ 厘米}\)，由勾股定理得
\[
(5+x)^2+\left(\dfrac{x}{2}\right)^2=10^2
\]
展开并化简，得到 \(x^2+8x-60=0\). 解得
\[
x=\dfrac{-8\pm\sqrt{8^2+240}}{2}=\dfrac{-8\pm4\sqrt{19}}{2}=-4\pm2\sqrt{19}
\]
边长必须为正，所以舍去 \(-4-2\sqrt{19}\)，取 \(x=2\sqrt{19}-4\text{ 厘米}\). 数值上 \(2\sqrt{19}-4\approx4.72\)，故正方形的边长约为 \(4.7\text{ 厘米}\).
\end{solution}
